\section{Scope, conventions, and the results established}
\label{intro:section}

This article continues the exact-identity programme in the unified
ProveIt manuscript and the six archives present in its incoming
directory at commit
\begin{center}\sourcepin.\end{center}
The commit is dated 11 October 2026, 18:22:01 UTC. The incoming archive
interiors were retrieved and inspected for overlap with the selected
targets. The repository manuscript and the relevant primary references
were checked where the proofs use them. This is a targeted mathematical
audit of those dependencies, not an exhaustive certification of every
statement in the repository. The deliverable is a separate article;
repository integration can preserve the existing thematic organization.

The main technical theme is to retain exact finite differences and
finite asymptotic terms before taking a limit. For tails, the missing
information is an endpoint product. For zeta continuation, it is a
finite residue contribution. For moment reconstruction, it is an exact
rational determinant rather than a floating rank test. These three
mechanisms lead to complementary identities among the functions named
in the research programme.

\subsection{A map from source questions to the present results}

\begin{longtable}{@{}>{\raggedright\arraybackslash}p{0.24\textwidth}>{\raggedright\arraybackslash}p{0.46\textwidth}>{\raggedright\arraybackslash}p{0.24\textwidth}@{}}
\toprule
Source target & Result established here & Precise status\\
\midrule
\endfirsthead
\toprule Source target & Result established here & Precise status\\\midrule
\endhead
Harmonic-tail square identities & A rational-tail finite certificate,
continuous order, all-index Stieltjes products, and anchored Gamma
integration (Section~\ref{ats:section}) & Additional identities with
ordinary convergence.\\[5pt]
Shifts and Twists, question Q6 \cite{HTST} & A finite generator for
arbitrary root-of-unity inner and outer colors
(Theorem~\ref{cyc:master}) & Solved for the specified parity/conjugation
sector; the complementary sector remains open.\\[5pt]
Quadratic-Spectrum Stieltjes, question 2 \cite{QSpec} & Every order
$\Delta_r$ for every polynomial weight
(Theorem~\ref{sd:thm:all-orders}) & Complete explicit correction formula;
generic finite polylogarithm reduction is not claimed.\\[5pt]
The same report, question 7 & Any number of quadratic factors, every
even polynomial weight, and every coincident-mass coefficient
(Theorems~\ref{mm:thm:second}--\ref{mm:thm:confluence}) & Complete at the
second spectral derivative in this setting.\\[5pt]
Optimal moment prefix, order eight \cite{PRR} & Full formal rank for
$p=8$, $q=2b$, $b\ge5$, outside $8001\pmod{16384}$
(Theorem~\ref{rt:rank8}) & Strengthens the existing exclusion
$1\pmod{16}$; one smaller progression remains.\\[5pt]
Optimal moment prefix, order ten & Full formal rank for $p=10$,
$q=2b$, $b\ge6$, outside $962\pmod{2048}$
(Theorem~\ref{rt:rank10}) & New infinite family and exact finite base
cases; no all-order theorem is inferred.\\
\bottomrule
\end{longtable}

The gap-four moment-prefix theorem is already established in the two
latest harmonic-tail reports \cite{HTGC,HTST}. It is not reproved here
as a new result. The Gaussian weight-six and revised weight-eight
numerical conjectures are not settled by this article. The new
cyclotomic identities concern different, explicitly specified sums;
their existence is not a missing proof of those numerical equalities.

\subsection{Conventions fixed throughout}

The Hurwitz zeta function is initially
$\zeta(s,a)=\sum_{n\ge0}(n+a)^{-s}$ for $\Re s>1$, $\Re a>0$,
then continued in its usual variables. Our generalized Stieltjes
constants are defined by
\begin{equation}\label{intro:stieltjes}
 \zeta(1+u,a)=\frac1u+
       \sum_{r\ge0}\frac{(-1)^r}{r!}\gamma_r(a)u^r.
\end{equation}
Thus $\gamma_0(a)=-\psi(a)$, $\gamma_0(1)=\gamma$, and a literal
linear Laurent coefficient is $-\gamma_1(a)$.
The prime in $\zeta'(s,a)$ always differentiates $s$; a derivative
in $a$ is written $\partial_a$. We use
\[
 \psi(a)=\frac{\Gamma'(a)}{\Gamma(a)},\qquad
 \psi^{(r)}(a)=(-1)^{r+1}r!\zeta(r+1,a)\quad(r\ge1),
\]
and $H_n^{(r)}=\sum_{k=1}^{n}k^{-r}$, with $H_n=H_n^{(1)}$ and
$H_0=0$. The polylogarithm is
$\operatorname{Li}_s(z)=\sum_{k\ge1}z^k/k^s$ in its convergence
disk, with continuation or ordinary boundary values stated when used.
These conventions agree with the standard references
\cite{DLMFZeta,DLMFPolylog,DLMFPolygamma}.

All complex powers use principal logarithms on the initial positive
shift domains. Any later continuation or antiderivative specifies its
domain and its anchor. A square of a complex tail means its algebraic
square; a modulus square is written explicitly. A series called
\emph{ordinary} is a limit of its stated partial sums. A continued
spectral value is never silently substituted for a divergent ordinary
sum. ``Locally normally convergent'' means locally uniform absolute
convergence, sufficient for holomorphic parameter differentiation.

The symbols used for colors in Section~\ref{cyc:sec} are roots of
unity; the positive real shift and mass variables in
Sections~\ref{sd:section}--\ref{mm:section} are local to those sections.
Nielsen's alternating Hurwitz function $\mathcal B_s$ is distinguished
from the Bernoulli polynomials $B_j$ and the cyclotomic generating
function $B_\omega$. The moment-rank statements concern a formal
coordinate space in which only the specified exact reduction relations
have been imposed. They make no arithmetic independence assertion
about actual zeta values.

\subsection{What counts as proof in the companion package}

The infinite-dimensional assertions below have written analytic proofs,
including convergence, pole cancellation, and branch control. Exact
symbolic certificates verify the finite algebra used in those proofs.
The moment nonvanishing theorems combine a uniform Bernoulli congruence
with exact polynomial identities and finitely many rational base cases.
High-precision evaluations supply independent diagnostics for signs,
factors, and endpoint conventions. They are not interval certificates,
proofs of an infinite family, or proofs of period independence.
No statement in this delivery is claimed to have been formalized in
Lean or another proof assistant.

